\begin{figure}[!htbp]
\centering
\resizebox{0.96\textwidth}{!}{%
\begin{tikzpicture}[
  part/.style={draw=gray!55, fill=gray!8, rounded corners=3pt,
    minimum width=2.08cm, minimum height=0.58cm, align=center, font=\scriptsize\bfseries},
  msg/.style={->, >=Stealth, thick, accent!85, shorten <=2pt, shorten >=2pt},
  ret/.style={->, >=Stealth, dashed, gray!70, thick, shorten <=2pt, shorten >=2pt},
  life/.style={gray!28, dashed},
  lab/.style={font=\tiny, fill=white, inner sep=1.2pt, align=center}
]

\foreach \x/\name/\id in {
  0/Operator/op,
  2.45/Console/ui,
  4.9/FastAPI + job/api,
  7.35/Orchestrator/graph,
  9.8/Specialist/agent,
  12.25/MCP + Gemini/tools,
  14.7/Postgres/db
}{
  \node[part] (\id) at (\x,0) {\name};
  \draw[life] (\x,-0.42) -- (\x,-8.85);
}

\draw[msg] (op.south) ++(0,-0.45) -- node[lab, above]{submit URL} ++(2.45,0);
\draw[msg] (ui.south) ++(0,-0.78) -- node[lab, above]{POST workflow run} ++(2.45,0);
\draw[msg] (api.south) ++(0,-1.48) -- node[lab, above, text width=2.2cm]{create job and trace} ++(9.8,0);
\draw[ret] (api.south) ++(0,-2.18) -- node[lab, above]{run id} ++(-2.45,0);
\draw[msg] (api.south) ++(0,-2.88) -- node[lab, above]{run\_pipeline(url)} ++(2.45,0);
\draw[msg] (graph.south) ++(0,-3.58) -- node[lab, above, text width=2.25cm]{classification and handoffs} ++(2.45,0);
\draw[msg] (agent.south) ++(0,-4.28) -- node[lab, above, text width=2.1cm]{model turns and browser tools} ++(2.45,0);
\draw[ret] (tools.south) ++(0,-4.98) -- node[lab, above, text width=2.25cm]{observations and model output} ++(-2.45,0);
\draw[ret] (agent.south) ++(0,-5.68) -- node[lab, above]{typed result} ++(-2.45,0);
\draw[msg] (graph.south) ++(0,-6.38) -- node[lab, above, text width=2.35cm]{persist result and telemetry} ++(7.35,0);
\draw[msg] (ui.south) ++(0,-7.08) -- node[lab, above]{GET run detail} ++(2.45,0);
\draw[ret] (api.south) ++(0,-7.78) -- node[lab, above, text width=2.2cm]{evidence and trace payload} ++(-2.45,0);
\draw[ret] (ui.south) ++(0,-8.48) -- node[lab, above]{review} ++(-2.45,0);

\end{tikzpicture}
}
\caption{Sequence diagram for launching, executing, persisting, and reviewing a workflow run.}
\label{fig:run-sequence}
\end{figure}
